Der Satz von Cantor--Schröder--Bernstein besagt: Gibt es Injektionen
von \(A\) nach \(B\) und von \(B\) nach \(A\), so gibt es eine
Bijektion zwischen beiden Mengen. Umgekehrt liefern eine Bijektion und
ihre Umkehrfunktion sofort Injektionen in beide Richtungen.

\CSBHeading{Die Idee der Konstruktion}
Seien \(F\colon A\inj B\) und \(G\colon B\inj A\) die gegebenen
Injektionen. Wir suchen eine Teilmenge \(C\subseteq A\), für die folgende
beiden Einschränkungen Bijektionen sind:
\[
  \begin{array}{c@{\qquad}c@{\qquad}c}
    \text{Teilmenge von }A && \text{Teilmenge von }B\\[4pt]
    C & \xrightarrow{\ F\ } & F[C]\\[5pt]
    A\setminus C & \xleftarrow{\ G\ } & B\setminus F[C].
  \end{array}
\]
Die obere Zuordnung benutzt \(F\). Die untere Zuordnung kehren wir um:
Jedem \(a\in A\setminus C\) soll das eindeutige
\(b\in B\setminus F[C]\) mit \(G(b)=a\) zugeordnet werden.
So werden die beiden disjunkten Teile von \(A\) bijektiv auf die beiden
disjunkten Teile von \(B\) abgebildet. Zusammen ergeben sie die gesuchte
Bijektion zwischen \(A\) und \(B\).

Die obere Zuordnung ist für jede Teilmenge \(C\subseteq A\) bijektiv
auf \(F[C]\), weil \(F\) injektiv ist. Bei der unteren Zuordnung
ist die Injektivität ebenfalls gegeben. Entscheidend ist also die Gleichung
\[
  G[B\setminus F[C]]=A\setminus C.
\]
Sie stellt sicher, dass jeder Punkt des verbleibenden Teils von \(A\)
genau ein Urbild im verbleibenden Teil von \(B\) besitzt.

Der folgende Beweis konstruiert dafür die kleinste Teilmenge
\(C\subseteq A\), die \(A\setminus G[B]\) enthält und für die
\(G[F[C]]\subseteq C\) gilt. Aus ihrer Minimalität folgt
\[
  C=(A\setminus G[B])\cup G[F[C]].
\]
Diese Fixpunktgleichung liefert die benötigte Bildgleichung und damit
die beiden Bijektionen des Schemas.
